\section*{第 \ref{ch:g11:mutations} 章——突变与遗传变异}

\begin{solution}{exo:g11:mutations:1}
\omterm{def:g10:universal-dna:information}{DNA} 中\omterm{def:g10:universal-dna:nucleotide}{核苷酸}序列发生的改变，
并传递给该\omterm{def:g10:cells-common-unit:cell}{细胞}的后代。替换、插入、缺失。
\end{solution}

\begin{solution}{exo:g11:mutations:2}
\texttt{ATGCCAGAAT}：替换（第 6 位由 T 变为 A）。
\texttt{ATGCCTGAT}：缺失一个 A。\texttt{ATGCCTTGAAT}：
插入一个 T。
\end{solution}

\begin{solution}{exo:g11:mutations:3}
体\omterm{def:g10:cells-common-unit:cell}{细胞}\omterm{def:g11:mutations:mutation}{突变}：发生于体\omterm{def:g10:cells-common-unit:cell}{细胞}，仅传给该个体内
该\omterm{def:g10:cells-common-unit:cell}{细胞}的后代。生殖\omterm{def:g10:cells-common-unit:cell}{细胞}\omterm{def:g11:mutations:mutation}{突变}：发生于配子或其前体，
传给子代及其全部\omterm{def:g10:cells-common-unit:cell}{细胞}。只有生殖\omterm{def:g10:cells-common-unit:cell}{细胞}\omterm{def:g11:mutations:mutation}{突变}
才能进入遗传。
\end{solution}

\begin{solution}{exo:g11:mutations:4}
紫外线、X 射线、烟草烟雾中的苯并芘（还有
放射性、黄曲霉毒素）。紫外线使同一链上相邻的两个胸腺嘧啶
连接在一起，使\omterm{prop:g10:universal-dna:helix}{双螺旋}结构变形。
\end{solution}

\begin{solution}{exo:g11:mutations:5}
完整的互补链。因为每条链都决定
另一条链的序列，所以受损片段的正确序列
可以依据其互补链重建。
\end{solution}

\begin{solution}{exo:g11:mutations:6}
正常\omterm{def:g10:cells-common-unit:cell}{细胞}：约 6\%；修复缺陷型：0.01\%。修复使存活率
在该剂量下提高约 600 倍。
\end{solution}

\begin{solution}{exo:g11:mutations:7}
$20/10^9 = \num{2e-8}$。否：平板只显示哪些\omterm{def:g10:cells-common-unit:cell}{细胞}
原本就具有抗性；抗生素是否引发了\omterm{def:g11:mutations:mutation}{突变}，需要
通过影印检验来判断。
\end{solution}

\begin{solution}{exo:g11:mutations:8}
每个影印平板都接种自同一批原始菌落。
若抗生素是在接触时才诱发抗性，则受影响的\omterm{def:g10:cells-common-unit:cell}{细胞}
会随机分布在各平板的少数位置，且位置互不相同。
若抗性总是出现在相同位置，则说明这些原始菌落
在未用药的情况下培养时，本身就已含有抗性\omterm{def:g10:cells-common-unit:cell}{细胞}。
\end{solution}

\begin{solution}{exo:g11:mutations:9}
每个\omterm{def:g10:universal-dna:gene}{基因}拷贝的误差数：$2 \times 3000 \times 10^{-9} =
\num{6e-6}$（两条链）。在最后一次分裂产生的 $10^{12}$ 个\omterm{def:g10:cells-common-unit:cell}{细胞}中，
约 $\num{6e6}$ 个携带该\omterm{def:g10:universal-dna:gene}{基因}的新\omterm{def:g11:mutations:mutation}{突变}。
\end{solution}

\begin{solution}{exo:g11:mutations:10}
芽\omterm{def:g10:cells-common-unit:cell}{细胞}中发生的一次体\omterm{def:g10:cells-common-unit:cell}{细胞}\omterm{def:g11:mutations:mutation}{突变}使该枝条呈现这种颜色；扦插
材料正是该组织，因此保留该性状。花的生殖\omterm{def:g10:cells-common-unit:cell}{细胞}来自
同一枝条，但一个\omterm{def:g10:universal-dna:gene}{等位基因}的改变只传给
一半配子；胚珠则由携带正常\omterm{def:g10:universal-dna:gene}{等位基因}的
花粉受精。若新\omterm{def:g10:universal-dna:gene}{等位基因}为隐性，则会在
幼苗中被掩盖。
\end{solution}

\begin{solution}{exo:g11:mutations:11}
缺失的酶负责修复紫外线损伤；紫外线只能穿透
皮肤。肠道不受紫外线照射，也不会发生此类
损伤，因而不需要该修复功能。
\end{solution}

\begin{solution}{exo:g11:mutations:12}
损伤与剂量成正比：紫外线减少为十分之一，意味着
每个\omterm{def:g10:cells-common-unit:cell}{细胞}每小时的损伤也减少为十分之一——仍有数千处。晒伤
意味着大量\omterm{def:g10:cells-common-unit:cell}{细胞}死亡；低于该剂量时，\omterm{def:g10:cells-common-unit:cell}{细胞}仍可存活，但
未被修复的损伤会转变为\omterm{def:g11:mutations:mutation}{突变}。没有晒伤
只说明损伤可以承受，并非不存在。
\end{solution}

\begin{solution}{exo:g11:mutations:13}
本身都不是。在疟疾流行的地区，单个拷贝
可提高生存率，因而该\omterm{def:g10:universal-dna:gene}{等位基因}在种群中得以保留，尽管纯合携带者
患病；在无疟疾的地区，它只会致病，因而较为罕见。
环境和拷贝数共同决定一个
\omterm{def:g11:mutations:mutation}{突变}的价值。
\end{solution}

\begin{solution}{exo:g11:mutations:14}
敏感\omterm{def:g10:cells-common-unit:cell}{细胞}死亡；原本就存在的抗性\omterm{def:g10:cells-common-unit:cell}{细胞}
自由分裂，一天之内培养物便完全具有抗性。疗程
若中途停止，会留下最后的存活者——抗性最强的个体——
使其重新生长，体内抗性\omterm{def:g10:cells-common-unit:cell}{细胞}的比例随之升高；
坚持服完整个疗程，则可在它们繁殖之前将其杀灭。
\end{solution}

\begin{solution}{exo:g11:mutations:15}
大多数\omterm{def:g11:mutations:mutation}{突变}是中性的（位于\omterm{def:g10:universal-dna:gene}{基因}之外，或对
\omterm{def:g10:chemistry-of-life:families}{蛋白质}无影响），少数有害，极少数有益；\omterm{def:g11:mutations:mutation}{突变}是随机的，
不论是否有用都会发生，但选择会保留有益\omterm{def:g11:mutations:mutation}{突变}，
淘汰有害\omterm{def:g11:mutations:mutation}{突变}。有益与否取决于环境，因此某一\omterm{def:g11:mutations:mutation}{突变}
今天有害，在条件改变时可能恰恰变得重要。
经过数千代、数十亿个个体的积累，
即使罕见的有益\omterm{def:g11:mutations:mutation}{突变}也会反复出现：演化依托
经过筛选保留下来的\omterm{def:g11:mutations:mutation}{突变}，而非平均\omterm{def:g11:mutations:mutation}{突变}。
\end{solution}

\begin{solution}{pb:g11:mutations:1}
\textbf{1.} 使每个菌落都是来源已知的克隆，并使
菌落之间相隔足够远，以便在影印平板上可以区分。

\textbf{2.} 取自每个菌落顶部的几个\omterm{def:g10:cells-common-unit:cell}{细胞}，
其几何排列与平板上的相同，
因为绒布是平压着按下去、不滑动的。

\textbf{3.} 抗性\omterm{def:g10:cells-common-unit:cell}{细胞}已存在于那两个原始菌落中：
三个独立的复制平板不会偶然得到一致的结果。

\textbf{4.} 三块平板上的抗性菌落位于随机、彼此无关的位置，
因为哪些\omterm{def:g10:cells-common-unit:cell}{细胞}被诱导，
在各复制平板之间会各不相同。

\textbf{5.} \omterm{def:g11:mutations:mutation}{突变}发生在接触之前：
抗性与原始菌落相对应，而这些菌落从未接触过药物。

\textbf{6.} 来自两个抗性位置的\omterm{def:g10:cells-common-unit:cell}{细胞}能在链霉素上生长；
来自其他位置的\omterm{def:g10:cells-common-unit:cell}{细胞}则不能。

\textbf{7.} 大概只有部分\omterm{def:g10:cells-common-unit:cell}{细胞}如此：\omterm{def:g11:mutations:mutation}{突变}发生在菌落生长过程中的某次分裂，
只有该\omterm{def:g10:cells-common-unit:cell}{细胞}的后代具有抗性。
要判断这一点，可以将已知数量的菌落\omterm{def:g10:cells-common-unit:cell}{细胞}涂布于链霉素平板上，
并计数存活者。

\textbf{8.} 在 23 次分裂中的第 20 次发生：$1/2^{20-1}$\dots 更简单地说，
\omterm{def:g11:mutations:mutation}{突变}\omterm{def:g10:cells-common-unit:cell}{细胞}随后再分裂 3 次，得到 $2^3 = 8$ 个抗性\omterm{def:g10:cells-common-unit:cell}{细胞}，
在 $2^{23}$ 个\omterm{def:g10:cells-common-unit:cell}{细胞}中约占百万分之一。在第 3 次分裂时发生：
此时存在的 8 个\omterm{def:g10:cells-common-unit:cell}{细胞}之一，因此占菌落的八分之一。

\textbf{9.} 即使绒布取下的约 $10^4$ 个\omterm{def:g10:cells-common-unit:cell}{细胞}中只有几个抗性\omterm{def:g10:cells-common-unit:cell}{细胞}，
也足以在抗生素平板上形成菌落；
该检测的是抗性\omterm{def:g10:cells-common-unit:cell}{细胞}是否存在，而非其比例。

\textbf{10.} 该论证依赖于位置：
只有几何排列得以保持，
复制平板上的菌落才能与原始平板上的菌落相对应。

\textbf{11.} $30/\num{10000} = \num{3e-3}$.

\textbf{12.} 约 $\num{3e-3}$ 次\omterm{def:g11:mutations:mutation}{突变}出现在 $10^7$ 次分裂中：
即每次分裂 $\num{3e-10}$。

\textbf{13.} 与单个\omterm{def:g10:universal-dna:nucleotide}{核苷酸}的错误率属于同一数量级：
抗性\omterm{def:g11:mutations:mutation}{突变}本质上是一种特定的替换，
逃过了修复。

\textbf{14.} $10^9$ 次分裂、每次 $\num{3e-10}$：
平均而言仅相当于一次\omterm{def:g11:mutations:mutation}{突变}的一部分，因此多数培养物中没有抗性\omterm{def:g10:cells-common-unit:cell}{细胞}，或只有少数
（每个\omterm{def:g11:mutations:mutation}{突变}之后会随其后的分裂而倍增）。

\textbf{15.} 早期分裂中发生的\omterm{def:g11:mutations:mutation}{突变}会被培养物中的大部分\omterm{def:g10:cells-common-unit:cell}{细胞}继承，
因此计数较大；较晚发生的\omterm{def:g11:mutations:mutation}{突变}则只产生少数抗性\omterm{def:g10:cells-common-unit:cell}{细胞}。
由于发生的时间是随机的，
平行培养物的计数差异很大——这本身说明\omterm{def:g11:mutations:mutation}{突变}是在接触抗生素之前发生的，
而非为应对抗生素而产生。

\textbf{16.} 稀有：\omterm{def:g11:mutations:mutation}{突变}在每次分裂中很少发生。
随机：\omterm{def:g11:mutations:mutation}{突变}的发生与其是否有用无关。
该实验确立了随机性（位置论证）；
第三部分中的计数则用于测定稀有程度。

\textbf{17.} 它对种群进行了筛选：
杀死了敏感\omterm{def:g10:cells-common-unit:cell}{细胞}，使抗性\omterm{def:g10:cells-common-unit:cell}{细胞}得以增殖——这就是选择。

\textbf{18.} \omterm{def:g11:mutations:mutation}{突变}在任何需要出现之前，随机地提供了抗性\omterm{def:g10:universal-dna:gene}{等位基因}；
抗生素对种群进行了筛选，
只保留了其携带者。

\textbf{19.} \omterm{def:g11:mutations:mutation}{突变}改变的是参与药物作用的分子，
而非在没有药物时菌落的生长或外观；
\omterm{def:g10:universal-dna:gene}{基因}差异只有在它起作用的环境中，
才会产生可见的效应。

\textbf{20.} 三个复制平板上相同的位置证明，
抗性\omterm{def:g11:mutations:mutation}{突变}发生于接触抗生素之前，且与之无关；
测得的速率约为每个\omterm{def:g10:cells-common-unit:cell}{细胞}每次分裂
$\num{3e-10}$。
\end{solution}
